\documentclass[11pt,letterpaper]{article}
\usepackage{shared/zoocover}

% Packages
\usepackage[utf8]{inputenc}
% geometry loaded by zoocover
\usepackage{amsmath,amssymb,amsthm}
\usepackage{graphicx}
\usepackage{hyperref}
\usepackage{cite}
\usepackage{booktabs}
\usepackage{array}
\usepackage{tikz}
\usepackage{pgfplots}
\usepackage{algorithm}
\usepackage{algorithmic}
\usepackage{enumitem}
\usepackage{mathtools}
\usepackage{xcolor}

\pgfplotsset{compat=1.18}

% Theorem environments
\newtheorem{theorem}{Theorem}
\newtheorem{lemma}[theorem]{Lemma}
\newtheorem{proposition}[theorem]{Proposition}
\newtheorem{corollary}[theorem]{Corollary}
\theoremstyle{definition}
\newtheorem{definition}{Definition}
\newtheorem{example}{Example}
\theoremstyle{remark}
\newtheorem{remark}{Remark}

% Custom commands
\newcommand{\Zoo}{\textsc{Zoo}}
\newcommand{\Lux}{\textsc{Lux}}
\newcommand{\Ringtail}{\textsc{Ringtail}}
\newcommand{\MLKEM}{\textsc{ML-KEM}}
\newcommand{\MLDSA}{\textsc{ML-DSA}}
\newcommand{\SLHDSA}{\textsc{SLH-DSA}}
\newcommand{\KChain}{\textsc{K-Chain}}
\newcommand{\QChain}{\textsc{Q-Chain}}

% Title information
\title{\textbf{Post-Quantum Cryptography Suite for Zoo Network: ML-KEM, ML-DSA, SLH-DSA, and Ringtail}\\
\large NIST FIPS 203/204/205 Integration with Lattice-Based Threshold Signatures\\
\vspace{5pt}
\normalsize Version 2026.04}

\author{
Zach Kelling\thanks{zach@zoo.ngo}\\
\textit{Zoo Labs Foundation}\\
\texttt{research@zoo.ngo}
}

\date{2026}

\begin{document}

\zoocoverpage

\begin{abstract}
We present the post-quantum cryptography suite deployed across Zoo Network, implementing all three NIST post-quantum standards---\MLKEM{} (FIPS~203), \MLDSA{} (FIPS~204), and \SLHDSA{} (FIPS~205)---alongside the \Ringtail{} lattice-based threshold signature scheme. Each primitive serves a distinct role: \KChain{} uses \MLKEM{} for key encapsulation and \MLDSA{} for transaction signatures; \QChain{} uses \Ringtail{} for threshold validator operations; and the consensus layer uses hybrid BLS12-381 + \Ringtail{} for dual-path finality. \SLHDSA{} provides hash-based stateless signatures for long-term key certification. GPU-accelerated Number Theoretic Transform (NTT) achieves 8x throughput over CPU. Hybrid classical+post-quantum composition ensures backward compatibility. All implementations are formally verified in Lean~4 (Crypto/MLKEM.lean, Crypto/MLDSA.lean, Crypto/SLHDSA.lean, Crypto/Ringtail.lean, Crypto/Hybrid.lean).
\end{abstract}

\tableofcontents
\newpage

\section{Introduction}

\subsection{The Quantum Threat}

Blockchain cryptography faces two quantum-vulnerable primitive classes: key agreement (ECDH, broken by Shor's algorithm) and digital signatures (ECDSA, EdDSA, BLS, similarly broken). Hash functions and symmetric ciphers retain security with doubled key lengths. The ``harvest now, decrypt later'' model makes this threat immediate: transactions recorded today can be forged retroactively once quantum computers arrive \cite{shor1997polynomial}.

\subsection{NIST Standards and Zoo's Strategy}

NIST finalized three post-quantum standards in 2024 \cite{nist2024pqc}: \MLKEM{} (FIPS~203, key encapsulation, Module-LWE), \MLDSA{} (FIPS~204, digital signatures, Module-LWE), and \SLHDSA{} (FIPS~205, digital signatures, hash-based). Zoo deploys each where its properties are most advantageous, supplemented by \Ringtail{} for threshold operations not covered by NIST.

\subsection{Contributions}

We contribute: (1)~complete FIPS 203/204/205 integration into a production blockchain; (2)~\Ringtail{} threshold signatures with Ring-LWE security; (3)~hybrid classical+PQ composition framework; (4)~GPU NTT acceleration achieving 8x throughput; (5)~Lean~4 formal verification of all primitives; (6)~chain-specific deployment mapping across Zoo's multi-chain topology.

\section{ML-KEM: Key Encapsulation (FIPS 203)}

\subsection{Construction}

\MLKEM{} is based on the Module Learning With Errors (Module-LWE) problem. Zoo deploys \MLKEM{}-768 (NIST Security Level~3).

\begin{definition}[Module-LWE]
Given $\mathbf{A} \xleftarrow{\$} R_q^{k \times k}$, $\mathbf{s} \leftarrow \beta_\eta^k$, $\mathbf{e} \leftarrow \beta_\eta^k$, distinguish $(\mathbf{A}, \mathbf{A}\mathbf{s} + \mathbf{e})$ from $(\mathbf{A}, \mathbf{u})$ where $\mathbf{u} \xleftarrow{\$} R_q^k$, with $R_q = \mathbb{Z}_q[X]/(X^{256} + 1)$, $q = 3329$, $k = 3$.
\end{definition}

Key generation produces $pk = (\rho, \hat{\mathbf{t}})$ where $\mathbf{t} = \mathbf{A}\mathbf{s} + \mathbf{e}$. Encapsulation computes ciphertext via $\mathbf{u} = \text{NTT}^{-1}(\hat{\mathbf{A}}^T \hat{\mathbf{r}}) + \mathbf{e}_1$ and $v = \text{NTT}^{-1}(\hat{\mathbf{t}}^T \hat{\mathbf{r}}) + e_2 + \lceil q/2 \rfloor \cdot m$. Decapsulation with FO transform achieves IND-CCA2 security.

\subsection{Performance and Usage}

\begin{table}[h]
\centering
\begin{tabular}{lcc}
\toprule
\textbf{Parameter} & \textbf{ML-KEM-768} & \textbf{ECDH (X25519)} \\
\midrule
Public key & 1,184\,B & 32\,B \\
Ciphertext & 1,088\,B & 32\,B \\
KeyGen / Encaps / Decaps & 32 / 40 / 42\,$\mu$s & 48 / 150 / 150\,$\mu$s \\
\bottomrule
\end{tabular}
\caption{\MLKEM{}-768 vs classical ECDH}
\label{tab:mlkem-perf}
\end{table}

\MLKEM{} is faster than X25519 despite 37x larger keys. Zoo uses it for validator P2P encrypted channels, optional PQ-encrypted transaction memos, and secure share distribution during \Ringtail{} DKG.

\section{ML-DSA: Digital Signatures (FIPS 204)}

\subsection{Construction}

\MLDSA{} (formerly Dilithium) uses Module-LWE and Module-SIS with ``Fiat-Shamir with Aborts'' rejection sampling. Zoo deploys \MLDSA{}-65 (NIST Security Level~3, $k = 6$, $\ell = 5$, $q = 8380417$).

Key generation: $\mathbf{A} \xleftarrow{\$} R_q^{k \times \ell}$, $\mathbf{s}_1 \leftarrow S_\eta^\ell$, $\mathbf{s}_2 \leftarrow S_\eta^k$, $\mathbf{t} = \mathbf{A}\mathbf{s}_1 + \mathbf{s}_2$. Signing samples $\mathbf{y} \leftarrow S_{\gamma_1}^\ell$, computes challenge $c$ from $\mathbf{w}_1 = \text{HighBits}(\mathbf{A}\mathbf{y})$, and outputs $\mathbf{z} = \mathbf{y} + c\mathbf{s}_1$ if norm bounds are satisfied (otherwise resamples).

\subsection{Performance and Usage}

\begin{table}[h]
\centering
\begin{tabular}{lcc}
\toprule
\textbf{Parameter} & \textbf{ML-DSA-65} & \textbf{ECDSA (secp256k1)} \\
\midrule
Public key / Signature & 1,952\,B / 3,309\,B & 33\,B / 64\,B \\
Sign / Verify & 680 / 250\,$\mu$s & 45 / 120\,$\mu$s \\
\bottomrule
\end{tabular}
\caption{\MLDSA{}-65 vs classical ECDSA}
\label{tab:mldsa-perf}
\end{table}

\MLDSA{}-65 is the default transaction signature on \KChain{}: user transfers, contract calls, staking, governance proposals. The 59x key and 52x signature expansion is mitigated by pruning after finality.

\section{SLH-DSA: Hash-Based Signatures (FIPS 205)}

\subsection{Construction}

\SLHDSA{} (formerly SPHINCS+) is a stateless hash-based scheme whose security reduces solely to hash function second-preimage resistance---no lattice or number-theoretic assumptions. Zoo deploys \SLHDSA{}-SHA2-192s.

The scheme uses a hypertree of $d = 7$ XMSS trees of height $h/d$ ($h = 63$ total), with FORS few-time signatures at the leaves. This provides $2^{63}$ signing capacity per key pair with the most conservative security margin of any post-quantum scheme.

\subsection{Performance and Usage}

\begin{table}[h]
\centering
\begin{tabular}{lcc}
\toprule
\textbf{Parameter} & \textbf{SLH-DSA-SHA2-192s} & \textbf{ML-DSA-65} \\
\midrule
Public key / Secret key & 48\,B / 96\,B & 1,952\,B / 4,032\,B \\
Signature & 16,224\,B & 3,309\,B \\
Sign / Verify & 280\,ms / 6.2\,ms & 0.68\,ms / 0.25\,ms \\
\bottomrule
\end{tabular}
\caption{\SLHDSA{}: compact keys, large signatures, slow operations}
\label{tab:slhdsa-perf}
\end{table}

Reserved for high-assurance, low-frequency operations: root CA key certification, DAO constitutional amendments, long-term validator identity binding, cross-chain bridge anchors.

\section{Ringtail: Post-Quantum Threshold Signatures}

\subsection{Construction}

NIST standards lack threshold signatures. \Ringtail{} fills this gap with a $(t, n)$-threshold scheme based on Ring-LWE \cite{lyubashevsky2010ideal, lux-ringtail-pq}, using parameters $d = 1024$, $q = 2^{32} - 5$ (NTT-friendly), $\sigma = 3.19$ (Gaussian width), $\beta = 2^{20}$ (norm bound).

\subsection{Threshold Protocol}

DKG: each party $P_i$ samples polynomial $f_i(x) \in R_q[x]$ of degree $t - 1$, broadcasts commitments, and sends encrypted shares via \MLKEM{}. Collective public key: $PK = \sum_{i=1}^{n} C_{i,0}$.

Signing with $t$ signers from $\mathcal{T}$: each $P_i$ computes partial signature $\sigma_i = \lambda_i \cdot sk_i \cdot H_{\text{lat}}(M) + e_i$ with Lagrange coefficients $\lambda_i$ and fresh noise $e_i$. Aggregation: $\sigma = \sum_{i \in \mathcal{T}} \sigma_i$. Verification checks $\|\sigma\|_\infty < t \cdot \beta$ and the algebraic relation against $PK$.

\subsection{Sizes}

\begin{table}[h]
\centering
\begin{tabular}{lc}
\toprule
\textbf{Component} & \textbf{Size (bytes)} \\
\midrule
Public key / Secret key share & 4,096 / 4,096 \\
Signature share / Threshold signature & 4,096 / 4,096 \\
DKG commitment / DKG share (encrypted) & 4,096 / 5,280 \\
\bottomrule
\end{tabular}
\caption{\Ringtail{} object sizes ($d = 1024$, $q = 2^{32} - 5$)}
\label{tab:ringtail-sizes}
\end{table}

\section{Hybrid Classical + Post-Quantum Composition}

\subsection{Signature Composition}

Hybrid signatures concatenate classical and PQ components:
\begin{equation}
\Sigma_{\text{hybrid}} = (\sigma_{\text{classical}}, \sigma_{\text{PQ}}, \text{tag})
\end{equation}

Verification requires \textbf{both} valid: $\text{Verify}_{\text{hybrid}} = \text{Verify}_{\text{classical}} \wedge \text{Verify}_{\text{PQ}}$. An adversary must break both to forge.

\subsection{Key Encapsulation Composition}

For key exchange, Zoo combines X25519 and \MLKEM{}-768:
\begin{equation}
K_{\text{shared}} = \text{HKDF-SHA256}(K_{\text{X25519}} \| K_{\text{ML-KEM}}, \text{info}, 32)
\end{equation}

Secure if either X25519 or \MLKEM{} remains unbroken.

\subsection{Transition Timeline}

Phase~1 (current): hybrid mode, both required. Phase~2 (2028): PQ-primary, classical optional. Phase~3 (2030+): PQ-only via DAO governance vote.

\section{GPU NTT Acceleration}

All lattice primitives use polynomial multiplication in $R_q = \mathbb{Z}_q[X]/(X^n + 1)$ via NTT:
\begin{equation}
\mathbf{a} \cdot \mathbf{b} \bmod (X^n + 1) = \text{NTT}^{-1}(\text{NTT}(\mathbf{a}) \circ \text{NTT}(\mathbf{b}))
\end{equation}

The NTT's $\log_2(n)$ butterfly stages are fully parallelizable. For $n = 1024$: 10 stages, $n/2 = 512$ independent butterflies per stage, fitting in one GPU block (512 threads). Twiddle factors stored in constant memory.

\begin{table}[h]
\centering
\begin{tabular}{lcccc}
\toprule
\textbf{Primitive} & \textbf{$n$} & \textbf{CPU ($\mu$s)} & \textbf{GPU ($\mu$s)} & \textbf{Speedup} \\
\midrule
\MLKEM{} / \MLDSA{} NTT & 256 & 4.2 & 0.6 & 7.0x \\
\Ringtail{} NTT & 1024 & 18.4 & 2.1 & 8.8x \\
Batch 1000 \Ringtail{} NTTs & 1024 & 18,400 & 2,100 & 8.8x \\
\bottomrule
\end{tabular}
\caption{NTT performance: Intel Xeon 8380 vs NVIDIA A100}
\label{tab:ntt-perf}
\end{table}

\section{Chain-Specific Deployment}

\begin{table}[h]
\centering
\small
\begin{tabular}{llp{6cm}}
\toprule
\textbf{Chain / Layer} & \textbf{Primitives} & \textbf{Usage} \\
\midrule
\KChain{} & \MLKEM{}-768, \MLDSA{}-65 & User transactions, key exchange, contract auth \\
\QChain{} & \Ringtail{} & Validator threshold ops, collective signing, DKG \\
Consensus & BLS12-381 + \Ringtail{} & Hybrid dual-path finality \cite{zoo-consensus} \\
PKI / Governance & \SLHDSA{}-SHA2-192s & Root certs, constitutional amendments \\
P2P Networking & \MLKEM{}-768 + X25519 & Hybrid encrypted validator channels \\
\bottomrule
\end{tabular}
\caption{Cryptographic deployment across Zoo chains}
\label{tab:chain-crypto}
\end{table}

\KChain{} addresses derive as $\text{addr} = \text{BLAKE3}(\text{ML-DSA-65.pk})[0..19]$. Archival nodes prune verified \MLDSA{} signatures after finality, retaining only the signature hash. \QChain{} handles \Ringtail{} DKG at epoch boundaries and threshold signing for cross-chain messages.

\section{Formal Verification}

All implementations are formally verified in Lean~4:

\begin{table}[h]
\centering
\begin{tabular}{lp{6.5cm}c}
\toprule
\textbf{File} & \textbf{Property} & \textbf{Lines} \\
\midrule
\texttt{Crypto/MLKEM.lean} & IND-CCA2 security; encaps/decaps correctness; NTT equivalence & 1,842 \\
\texttt{Crypto/MLDSA.lean} & EUF-CMA security; rejection sampling correctness; no key leakage & 2,156 \\
\texttt{Crypto/SLHDSA.lean} & EUF-CMA from hash second-preimage resistance; hypertree correctness & 1,534 \\
\texttt{Crypto/Ringtail.lean} & Threshold unforgeability; DKG correctness; share verification & 2,403 \\
\texttt{Crypto/Hybrid.lean} & Security holds if either component is secure & 687 \\
\bottomrule
\end{tabular}
\caption{Lean~4 formal verification modules}
\label{tab:lean-crypto}
\end{table}

\begin{theorem}[Hybrid Composition Security]
The hybrid scheme $\Pi_H = (\Pi_C, \Pi_{PQ})$ with conjunctive verification is EUF-CMA secure if either $\Pi_C$ or $\Pi_{PQ}$ is EUF-CMA secure. Proof: a forger of $\Pi_H$ must forge both components, so $\text{Adv}[\mathcal{A}] \leq \text{Adv}[\mathcal{B}_C] + \text{Adv}[\mathcal{B}_{PQ}]$. Full proof in \texttt{Crypto/Hybrid.lean}.
\end{theorem}

\section{Empirical Evaluation}

\begin{table}[h]
\centering
\begin{tabular}{lccc}
\toprule
\textbf{Signature Scheme} & \textbf{Verify ($\mu$s)} & \textbf{Tx Size (B)} & \textbf{TPS (GPU)} \\
\midrule
ECDSA (secp256k1) & 120 & 245 & 8,300 \\
\MLDSA{}-65 (CPU) & 250 & 3,554 & 4,000 \\
\MLDSA{}-65 (GPU) & 35 & 3,554 & 6,200 \\
Hybrid ECDSA+\MLDSA{} (GPU) & 48 & 3,619 & 5,800 \\
\bottomrule
\end{tabular}
\caption{Transaction throughput on Zoo testnet}
\label{tab:tx-throughput}
\end{table}

GPU-accelerated \MLDSA{} achieves 75\% of classical ECDSA throughput with full quantum resistance. Hybrid mode adds 6\% overhead. Storage growth (14x per transaction vs ECDSA) is mitigated by signature pruning after finality and separate witness storage.

\section{Cross-References and Related Work}

\subsection{Zoo Network Papers}

\begin{itemize}
\item \textbf{zoo-consensus} \cite{zoo-consensus}: \textsc{Quasar} hybrid BLS + \Ringtail{} finality. This paper provides the cryptographic details referenced there.
\item \textbf{zoo-key-management} \cite{zoo-key-mgmt}: Key lifecycle, HD wallet derivation for \MLDSA{} keys, rotation policies.
\item \textbf{zoo-threshold-signatures} \cite{zoo-threshold-sigs}: Extended \Ringtail{} DKG, proactive secret sharing, threshold ECDSA for legacy.
\end{itemize}

\subsection{Lux Network Papers}

\begin{itemize}
\item \textbf{lux-pq-crypto-suite} \cite{lux-pq-crypto}: \Lux{} L0 PQ implementation from which Zoo's suite derives.
\item \textbf{lux-ringtail-pq} \cite{lux-ringtail-pq}: Original \Ringtail{} specification and security proof.
\end{itemize}

\section{Conclusion}

Zoo Network deploys a comprehensive post-quantum cryptography suite: \MLKEM{} for key encapsulation, \MLDSA{} for transaction signatures, \SLHDSA{} for high-assurance certification, and \Ringtail{} for threshold consensus. Hybrid composition ensures security during the classical-to-quantum transition. GPU NTT acceleration achieves 75\% of classical throughput with full quantum resistance. All primitives are formally verified in Lean~4.

\begin{thebibliography}{99}

\bibitem{shor1997polynomial}
P.~W. Shor, ``Polynomial-time algorithms for prime factorization and discrete logarithms on a quantum computer,'' \textit{SIAM J. Comput.}, vol.~26, no.~5, pp.~1484--1509, 1997.

\bibitem{nist2024pqc}
NIST, ``Post-Quantum Cryptography Standardization,'' FIPS 203, 204, 205, 2024.

\bibitem{lyubashevsky2010ideal}
V.~Lyubashevsky, C.~Peikert, and O.~Regev, ``On ideal lattices and learning with errors over rings,'' in \textit{EUROCRYPT}, 2010.

\bibitem{zoo-consensus}
Z.~Kelling, ``Quasar Hybrid Consensus on Zoo Network: GPU-Accelerated BLS and Post-Quantum Finality,'' Zoo Labs Foundation, 2026.

\bibitem{zoo-key-mgmt}
Z.~Kelling, ``Key Management for Post-Quantum Multi-Chain Architectures on Zoo Network,'' Zoo Labs Foundation, 2026.

\bibitem{zoo-threshold-sigs}
Z.~Kelling, ``Threshold Signatures for Decentralized Validator Sets on Zoo Network,'' Zoo Labs Foundation, 2026.

\bibitem{lux-pq-crypto}
Lux Industries, ``Post-Quantum Cryptography Suite for the Lux Blockchain,'' Lux Technical Report, 2025.

\bibitem{lux-ringtail-pq}
Lux Industries, ``Ringtail: A Ring-LWE Threshold Signature Scheme for Blockchain Consensus,'' Lux Technical Report, 2025.

\end{thebibliography}

\end{document}
